% 教学构造，不是实验结果。复杂度 t 从 1 到 9：
% training(t)=.3 exp[-.4(t-1)]+.02；validation(t)=.13+.015(t-5)^2。
% 最低验证误差位于 t=5；仅对选定模型标出独立测试误差 .17。
\begin{tikzpicture}[foundation picture]
 \path[use as bounding box] (0,-9) rectangle (112,61);
 \begin{scope}[xshift=4mm]
 \draw[foundation axis] (10,10)--(99,10);
 \draw[foundation axis] (10,10)--(10,52);
 \node[rotate=90] at (1,32) {经验误差};
 \node at (54,-6) {模型复杂度};
 \foreach \error/\label in {0/0,.2/0.2,.4/0.4}{
  \draw[FigureStroke,line width=.9pt] (9,10+80*\error)--(10,10+80*\error);
  \node[anchor=east] at (8,10+80*\error) {$\label$};
 }
 \draw[FigureBlue,line width=1.1pt] (15,57)--(25,57);
 \fill[FigureBlue] (20,57) circle (.9);
 \node[anchor=west] at (27,57) {训练误差};
 \draw[FigureRed,line width=1.1pt,dashed] (59,57)--(69,57);
 \draw[FigureRed,fill=white,line width=1.1pt] (62.9,55.9) rectangle (65.1,58.1);
 \node[anchor=west] at (71,57) {验证误差};
 \draw[FigureBlue,line width=1.1pt,domain=1:9,samples=80]
  plot ({10+10*(\x-1)},{10+80*(.3*exp(-.4*(\x-1))+.02)});
 \draw[FigureRed,line width=1.1pt,dashed,domain=1:9,samples=80]
  plot ({10+10*(\x-1)},{10+80*(.13+.015*(\x-5)^2)});
 \foreach \t in {1,...,9}{
  \pgfmathsetmacro{\tx}{10+10*(\t-1)}
  \pgfmathsetmacro{\trainy}{10+80*(.3*exp(-.4*(\t-1))+.02)}
  \pgfmathsetmacro{\valy}{10+80*(.13+.015*(\t-5)^2)}
  \fill[FigureBlue] (\tx,\trainy) circle (.8);
  \draw[FigureRed,fill=white,line width=1pt] (\tx-.85,\valy-.85) rectangle (\tx+.85,\valy+.85);
 }
 \draw[FigureMuted,line width=.9pt,dotted] (50,10)--(50,20.4);
 \draw[FigureRed,line width=1.1pt] (50,20.4) circle (2.2);
 \node[below] at (50,9) {$\hat k$};
 % 测试误差只在最终锁定的复杂度处出现；不绘制测试曲线。
 \fill[FigureInk] (50,25.2)--(48.7,22.8)--(51.3,22.8)--cycle;
 \node[anchor=center] (testerror) at (77,46) {测试误差};
 \draw[foundation flow,FigureMuted,line width=.9pt] (51.7,24.5)--(testerror.south);
 \end{scope}
\end{tikzpicture}%
